% !TEX root = ../../../main.tex
\section{Change of basis formula for vectors}
\label{sec:linear-algebra-i:change-of-basis-for-vectors}

The coordinates introduced by a basis describe a vector relative to that
basis. A second basis gives a second description of the same vector.
The matrix in this section translates between those descriptions. When a
linear transformation is applied, it produces an output vector,
as in the shear of
\autoref{fig:linear-algebra-i:shear-grid-overview}. Both calculations use
matrices because both reconstruct vectors from linear combinations.

Let \(E=[e_1,\ldots,e_n]\) and \(F=[f_1,\ldots,f_n]\) be two ordered
bases of \(V\). Each vector \(f_j\) has coordinates in the basis \(E\),
so there is a matrix \(M\in\mathbb F^{n\times n}\) whose \(j\)-th column
contains those coordinates. In matrix notation,
\begin{marginfigure}
  \centering
  % Same vector in two bases. For alignment only, the drawing uses
% e1=(1,0), e2=(0,1), f1=(1,1/2), f2=(-1/2,1), b=(2,7/5).
\begin{tikzpicture}[
  x=8mm,y=8mm,font=\footnotesize,line width=0.45pt,
  vector/.style={-{Latex[length=1.6mm,width=1.1mm]}},
  axis/.style={vector,black!40,line width=0.35pt},
  guide/.style={dashed,black!45,line width=0.35pt}
]
  \begin{scope}
    \coordinate (O) at (0,0);
    \coordinate (P) at (1.3,2.4);
    \draw[axis] (-0.4,0) -- (2.8,0);
    \draw[axis] (0,-0.3) -- (0,2.8);
    \draw[guide] (1.3,0) -- (P) -- (0,2.4);
    \draw[vector] (O) -- (1,0) node[midway,below=3pt] {\(e_1\)};
    \draw[vector] (O) -- (0,1) node[midway,left=3pt] {\(e_2\)};
    \draw[vector] (O) -- (P) node[above right=2pt] {\(v\)};
    \node[below=3pt] at (1.3,0) {\(a^1e_1\)};
    \node[left=3pt] at (0,2.4) {\(a^2e_2\)};
    \node at (1.1,3.4) {Basis \(E\)};
    \node at (1.1,-0.9) {\(v=a^1e_1+a^2e_2\)};
  \end{scope}
  \begin{scope}[yshift=-44mm]
    \coordinate (O) at (0,0);
    \coordinate (F1) at (1,0.5);
    \coordinate (F2) at (-0.5,1);
    \coordinate (C1) at ($2*(F1)$);
    \coordinate (C2) at ($1.4*(F2)$);
    \coordinate (P) at ($(C1)+(C2)$);
    \draw[axis] (-0.4,-0.2) -- (2.7,1.35);
    \draw[axis] (0.15,-0.3) -- (-1.25,2.5);
    \draw[guide] (C1) -- (P) -- (C2);
    \draw[vector] (O) -- (F1) node[below=3pt] {\(f_1\)};
    \draw[vector] (O) -- (F2) node[midway,left=3pt] {\(f_2\)};
    \draw[vector] (O) -- (P) node[above right=2pt] {\(v\)};
    \node[below right=2pt] at (C1) {\(b^1f_1\)};
    \node[left=3pt] at (C2) {\(b^2f_2\)};
    \node at (0.9,3.4) {Basis \(F\)};
    \node at (0.9,-0.9) {\(v=b^1f_1+b^2f_2\)};
  \end{scope}
  \path (current bounding box.south) ++(0,-2mm);
\end{tikzpicture}

  \caption[One vector in two bases]{The same vector \(v\) in two bases
    of the real plane, shown schematically. Dashed lines complete the
    components. The coordinates satisfy \(a=Mb\) when \(F=EM\).}
  \label{fig:linear-algebra-i:vector-basis-coordinates}
\end{marginfigure}
\[
  F=EM.
\]
Because \(F\) is also a basis, there is a matrix \(N\) with \(E=FN\).
Substituting these expressions into one another gives
\(E=EMN\) and \(F=FNM\). Uniqueness of coefficients in each basis
gives \(MN=I_n\) and \(NM=I_n\). Thus \(N=M^{-1}\). The change-of-basis
matrix is invertible, and it converts coordinates in either direction.

If \(v\in V\) has coordinate columns \(a\) in basis \(E\) and \(b\) in
basis \(F\), then \(v=Ea=Fb\). Using \(F=EM\),
\[
  v=Fb=EMb,
  \qquad a=Mb,
  \qquad b=M^{-1}a.
\]
Thus \(M\) converts \(F\)-coordinates to \(E\)-coordinates, and
\(M^{-1}\) converts \(E\)-coordinates to \(F\)-coordinates.

The vector stays fixed while its coordinate column changes.
\autoref{fig:linear-algebra-i:vector-basis-coordinates} shows the two
decompositions in the real plane.
